% =============================================================================
% INTRODUCTION — 5 minutes maximum
% =============================================================================
\section{Introduction}

% ~20 s
\begin{frame}[plain]
  \begin{tikzpicture}[remember picture,overlay]
    \fill[Navy] (current page.south west) rectangle (current page.north east);
    \fill[Blue] (current page.south west) rectangle ([yshift=0.055\paperheight]current page.south east);
  \end{tikzpicture}
  \vfill
  \begin{center}
    {\color{white}\Huge\bfseries Validation des blocs de compétences}\\[0.6em]
    {\color{white!88}\Large Expert en Architecture et Développement Logiciel}\\[1.5em]
    {\color{white}\large Thomas Graber}\\[0.3em]
    {\color{white!75}Learn IT · Brest Open Campus — Alternance au SHOM}
  \end{center}
  \vfill
\end{frame}

% ~50 s
\begin{frame}{Mon parcours et mon contexte professionnel}
\context{INTRO}{Contexte}
\begin{columns}[T]
\begin{column}{0.50\linewidth}
\begin{block}{Parcours}
\begin{tightitems}
  \item Licence informatique : systèmes, algorithmique, bases de données et développement.
  \item Mastère développement, programmation et Data/IA.
  \item Évolution progressive vers le \strong{back-end}, l'\strong{architecture logicielle} et la \strong{donnée}.
\end{tightitems}
\end{block}
\end{column}
\begin{column}{0.46\linewidth}
\begin{block}{Alternance au SHOM}
\begin{tightitems}
  \item Depuis septembre 2024, développement de services web et d'outils internes.
  \item Environnement sécurisé, métier et fortement contraint.
  \item Passage du pôle SYSPROD au pôle DATA.
  \item Travail sur l'analyse d'existant, le développement, la production et la documentation.
\end{tightitems}
\end{block}
\end{column}
\end{columns}
\end{frame}

% ~1 min 15
\begin{frame}{Huit projets, huit rôles complémentaires}
\context{INTRO}{Portefeuille}
\scriptsize
\begin{tabularx}{\linewidth}{>{\bfseries}p{3.05cm}p{2.05cm}X}
\toprule
Projet & Contexte & Rôle principal dans le dossier et l'oral \\
\midrule
FCU / Ediacara & Professionnel & Évolution d'une chaîne de production legacy, validation, exploitation et documentation. \\
S-123 & Professionnel & Architecture métier, méta-modèle S-100, XML/XSD, PostgreSQL et FastAPI. \\
FeastVerse & Académique & Back-end Spring Boot sécurisé, tests, documentation et déploiement. \\
Brise-glace & Académique & Front-end React/TypeScript, temps réel et intégration Supabase. \\
Prédiction des coûts d'assurance & Académique & Préparation des données, modèles prédictifs, explicabilité et passage à l'échelle. \\
Homelab & Personnel & GitOps, Kubernetes, CI/CD et observabilité. \\
piv2glz & Professionnel & Tests automatisés, intégration continue et construction de livrables Windows. \\
Lynx Immo & Académique / client réel & Faisabilité, planification, KPI, risques et pilotage d'équipe. \\
\bottomrule
\end{tabularx}
\end{frame}

% ~1 min 30
\begin{frame}{Une soutenance structurée par blocs, pas par projets}
\context{INTRO}{Carte de lecture}
\scriptsize
\begin{tabularx}{\linewidth}{>{\bfseries}p{2.95cm}ccccX}
\toprule
Projet & B1 & B2 & B3 & B4 & Angle principal à l'oral \\
\midrule
FCU / Ediacara & \checkmarkblue & \textbullet & \checkmarkblue & \textbullet & C1.3, C3.3, C3.6 \\
S-123 & \textbullet & \checkmarkblue &  &  & C2.1 \\
FeastVerse &  & \checkmarkblue & \textbullet &  & C2.4 \\
Brise-glace &  & \checkmarkblue &  &  & C2.3 \\
Prédiction assurance &  & \checkmarkblue &  &  & C2.5 \\
Homelab & \textbullet &  & \checkmarkblue &  & C3.4, C3.5 \\
piv2glz &  & \checkmarkblue & \checkmarkblue &  & C2.2, C3.1, C3.2 \\
Lynx Immo & \checkmarkblue &  &  & \checkmarkblue & C1.1--C1.2, C1.4--C1.7, C4.1--C4.6 \\
\bottomrule
\end{tabularx}
\vspace{0.55em}
\begin{block}{Principe de lecture}
Le dossier écrit raconte les projets. \textbf{La soutenance mobilise ces projets comme preuves des compétences du référentiel.}
\end{block}
\end{frame}
